\clearpage
\begin{proofbox}
\textbf{\textcolor{CProof}{方法一：中点与垂直法（门槛最低）}}
\begin{center}
\begin{minipage}{\linewidth}\centering
\begin{tikzpicture}[x=1.95cm,y=1.95cm,>=stealth,font=\small,line cap=round,line join=round]
\draw[->,gray!75] (-0.35,0)--(3.96,0) node[right=3pt]{$x$};
\draw[->,gray!75] (0,-0.35)--(0,3.96) node[above=3pt]{$y$};
\node at (-.3,.4) {$O$};
\draw[gray,dashed] (-0.15,-0.15)--(3.8,3.8) node[above left=3pt]{$y=x$};

\draw[gray!55,densely dotted] (1,0)--(1,3)--(0,3) (3,0)--(3,1)--(0,1);
\draw[teal,thick] (1,3)--(3,1);
\fill[teal] (1,3) circle(2pt) node[above left=5pt]{$P(1,3)$};
\fill[blue] (3,1) circle(2pt) node[below right=5pt]{$P'(3,1)$};
\fill (2,2) circle(1.8pt) node[right=12pt]{$M(2,2)$};
\draw (1.84,2.16)--(1.68,2)--(1.84,1.84);

\end{tikzpicture}\par\smallskip
{\small $PP^\prime\perp y=x,\qquad PM=MP^\prime$}\par
{\small 中点在镜线上，横纵坐标互换；两轴采用相同单位长度。}\par\smallskip
\end{minipage}
\end{center}


\begin{enumerate}[leftmargin=*,itemsep=.45em]
\item 任取 $P(a,b)$，猜想镜像为 $Q(b,a)$。
\item 中点 $M=\bigl((a+b)/2,(a+b)/2\bigr)$，两坐标相等，所以 $M\in y=x$。
\item 若 $a\ne b$，$k_{PQ}=(a-b)/(b-a)=-1$，而对称轴斜率为 $1$，故 $PQ\perp y=x$。
\item 对称轴垂直平分 $PQ$，所以两点关于 $y=x$ 对称。若 $a=b$，点在轴上，与自身重合。
\end{enumerate}
这证明交换坐标恰好就是轴对称，而不是只对某些特殊点成立。
\end{proofbox}
\clearpage
\begin{proofbox}
\textbf{\textcolor{CProof}{方法二：交换方程法（做题最快）}}
\begin{center}
\begin{minipage}{\linewidth}\centering
\begin{tikzpicture}[x=1.35cm,y=1.35cm,>=stealth,font=\small,line cap=round,line join=round]
\draw[->,gray!75] (-3.3,0)--(3.46,0) node[right=3pt]{$x$};
\draw[->,gray!75] (0,-3.3)--(0,3.46) node[above=3pt]{$y$};
\node at (-.3,.4) {$O$};
\draw[gray,dashed] (-3.3,-3.3)--(3.3,3.3) node[above left=3pt]{$y=x$};

\draw[teal,very thick,domain=-3:3,samples=121] plot({\x*\x/4},\x);
\draw[blue,very thick,domain=-3:3,samples=121] plot(\x,{\x*\x/4});
\draw[orange!80!black,densely dotted] (1,2)--(2,1);
\fill[teal] (1,2) circle(2pt) node[left=5pt]{$(1,2)$};
\fill[blue] (2,1) circle(2pt) node[below right=8pt]{$(2,1)$};
\draw[teal,dashed] (.3,1.3)--(1.7,2.7);
\draw[blue,dashed] (1.3,.3)--(2.7,1.7);

\end{tikzpicture}\par\smallskip
{\small \textcolor{teal}{\rule[.5ex]{1em}{.8pt}\ $y^2=4x$}\qquad\textcolor{blue}{\rule[.5ex]{1em}{.8pt}\ $x^2=4y$}}\par
{\small 对应切点 $(1,2)$ 与 $(2,1)$ 处，切线斜率均为 $1$。}\par\smallskip
\end{minipage}
\end{center}


\begin{enumerate}[leftmargin=*,itemsep=.45em]
\item 原曲线上任一点记为 $(u,v)$，满足 $F(u,v)=0$。
\item 镜像点记为 $(x,y)=(v,u)$，故 $u=y,v=x$。
\item 代回原式，得到 $F(y,x)=0$。
\item 反过来，若 $(x,y)$ 满足 $F(y,x)=0$，则 $(y,x)$ 满足原式。再交换一次回到 $(x,y)$，因此所得方程不多不少恰为完整镜像。
\end{enumerate}
若原方程是 $v=f(u)$，代入得 $x=f(y)$。如果每个 $x\in E$ 仅对应一个 $y\in D$，即可写为 $y=f^{-1}(x)$。否则只能保留关系式或分支。
\end{proofbox}
\clearpage
\begin{proofbox}
\textbf{\textcolor{CProof}{方法三：输入与输出互换（标准法）}}
\begin{center}
\begin{minipage}{\linewidth}\centering
\begin{tikzpicture}[x=1.35cm,y=1.35cm,>=stealth,font=\small,line cap=round,line join=round]
\draw[->,gray!75] (-3.3,0)--(3.46,0) node[right=3pt]{$x$};
\draw[->,gray!75] (0,-3.3)--(0,3.46) node[above=3pt]{$y$};
\node at (-.3,.4) {$O$};
\draw[gray,dashed] (-3.3,-3.3)--(3.3,3.3) node[above left=3pt]{$y=x$};

\draw[teal,very thick,domain=-3:1.65,samples=120] plot(\x,{2^\x});
\draw[blue,very thick,domain=-3:1.65,samples=120] plot({2^\x},\x);
\draw[orange!80!black,densely dotted] (1,2)--(2,1);
\fill[teal] (0,1) circle(2pt) node[left=4pt]{$(0,1)$};
\fill[blue] (1,0) circle(2pt) node[below right=8pt]{$(1,0)$};
\fill[teal] (1,2) circle(2pt) node[above left=4pt]{$(1,2)$};
\fill[blue] (2,1) circle(2pt) node[below right=4pt]{$(2,1)$};

\end{tikzpicture}\par\smallskip
{\small \textcolor{teal}{\rule[.5ex]{1em}{.8pt}\ $y=2^x$}\qquad\textcolor{blue}{\rule[.5ex]{1em}{.8pt}\ $y=\log_2x$}}\par
{\small 原函数与反函数的对应点交换横纵坐标。}\par\smallskip
\end{minipage}
\end{center}


设 $f:D\to E$ 一一对应，$E=f(D)$。
\begin{enumerate}[leftmargin=*,itemsep=.45em]
\item 任取 $a\in D$，令 $b=f(a)$，于是 $(a,b)$ 在原图上。
\item 按反函数定义，输入 $b$ 后返回唯一的 $a$，即 $f^{-1}(b)=a$。
\item 因而 $(b,a)$ 在反函数图上；反函数每一点也可这样还原到原图。
\item 所以二者关于 $y=x$ 对称；并有 $f^{-1}(f(a))=a$，以及 $f(f^{-1}(b))=b$。
\end{enumerate}
\textbf{单调性：}设 $u<v$ 都在 $E$，令 $s=f^{-1}(u),t=f^{-1}(v)$。若 $f$ 严格递增，则 $s\ge t$ 会推出 $u\ge v$，矛盾，所以 $s<t$；若 $f$ 严格递减，$s\le t$ 会推出 $u\ge v$，故 $s>t$。
\end{proofbox}
\clearpage
\begin{proofbox}
\textbf{\textcolor{CProof}{推论一：自身对称的充要条件}}
\begin{center}
\begin{minipage}{\linewidth}\centering
\begin{tikzpicture}[x=1.35cm,y=1.35cm,>=stealth,font=\small,line cap=round,line join=round]
\draw[->,gray!75] (-3.3,0)--(3.46,0) node[right=3pt]{$x$};
\draw[->,gray!75] (0,-3.3)--(0,3.46) node[above=3pt]{$y$};
\node at (-.3,.4) {$O$};
\draw[gray,dashed] (-3.3,-3.3)--(3.3,3.3) node[above left=3pt]{$y=x$};

\foreach \s in {-1,1}{\draw[teal,very thick,domain=.333333:3,samples=120] plot({\s*\x},{\s/\x});}
\draw[orange!80!black,densely dotted,thick] (2,.5)--(.5,2);
\fill[teal] (2,.5) circle(2pt) node[above right=6pt]{$(2,1/2)$};
\fill[teal] (.5,2) circle(2pt) node[above right=6pt]{$(1/2,2)$};

\end{tikzpicture}\par\smallskip
{\small \textcolor{teal}{$y=1/x\ (x\ne0)$：原函数与反函数图象完全重合}}\par
{\small 交换坐标后仍满足 $xy=1$；两条分支都必须保留。}\par\smallskip
\end{minipage}
\end{center}


\textbf{必要性：}若原图自身关于 $y=x$ 对称，任取 $x\in D$，点 $(x,f(x))$ 的镜像 $(f(x),x)$ 仍在原图。于是 $f(x)\in D$ 且 $f(f(x))=x$。

\textbf{充分性：}若 $f(D)\subseteq D$ 且任意 $x\in D$ 均满足 $f(f(x))=x$，则 $(f(x),x)$ 正是原图上一点。任一点镜像仍在原图；反射两次为原点，所以整图重合。

此外任取 $x\in D$，$x=f(f(x))\in f(D)$，所以 $D\subseteq f(D)$；结合原条件得 $f(D)=D$。若 $f(u)=f(v)$，两边再作用一次 $f$ 得 $u=v$，所以反函数存在且等于自身。
\end{proofbox}
\clearpage
\begin{proofbox}
\textbf{\textcolor{CProof}{推论二：严格递增函数的交点判定}}
\begin{center}
\begin{minipage}{\linewidth}\centering
\begin{tikzpicture}[x=1.35cm,y=1.35cm,>=stealth,font=\small,line cap=round,line join=round]
\draw[->,gray!75] (-3.3,0)--(3.46,0) node[right=3pt]{$x$};
\draw[->,gray!75] (0,-3.3)--(0,3.46) node[above=3pt]{$y$};
\node at (-.5,-.25) {$O$};
\draw[gray,dashed] (-3.3,-3.3)--(3.3,3.3) node[above left=3pt]{$y=x$};

\draw[teal,very thick] (-1,-3)--(2,3);
\draw[blue,very thick] (-3,-1)--(3,2);
\fill (1,1) circle(2.5pt) node[below right=8pt]{$(1,1)$};

\end{tikzpicture}\par\smallskip
{\small \textcolor{teal}{\rule[.5ex]{1em}{.8pt}\ $y=2x-1$}\qquad\textcolor{blue}{\rule[.5ex]{1em}{.8pt}\ $y=(x+1)/2$}}\par
{\small 严格递增时，公共点落在 $y=x$ 上。}\par\smallskip
\end{minipage}
\end{center}


若 $(a,b)$ 是原函数与反函数的交点，则 $b=f(a)$，且 $b=f^{-1}(a)$，后一式等价于 $a=f(b)$。因此 $a,b\in D$，并有
\[
 f(a)=b,\qquad f(b)=a.
\]
若 $a<b$，由严格递增得 $f(a)<f(b)$，代入即 $b<a$，矛盾；若 $a>b$ 同理矛盾。故 $a=b$。

反过来，若 $a\in D$ 且 $f(a)=a$，则 $a\in E$，$f^{-1}(a)=a$，所以 $(a,a)$ 为公共点。因此严格递增时求交点与解 $f(x)=x$ 完全等价。

若 $f$ 严格递增又自身对称，所有图上点都是交点，故 $f(x)=x$ 对所有 $x\in D$ 成立。
\end{proofbox}
\clearpage
\begin{proofbox}
\textbf{\textcolor{CProof}{推论三：切线斜率与反函数导数}}
\begin{center}
\begin{minipage}{\linewidth}\centering
\begin{tikzpicture}[x=1.35cm,y=1.35cm,>=stealth,font=\small,line cap=round,line join=round]
\draw[->,gray!75] (-3.3,0)--(3.46,0) node[right=3pt]{$x$};
\draw[->,gray!75] (0,-3.3)--(0,3.46) node[above=3pt]{$y$};
\node at (-.3,.4) {$O$};
\draw[gray,dashed] (-3.3,-3.3)--(3.3,3.3) node[above left=3pt]{$y=x$};

\draw[teal,very thick] (-1.5,-3)--(1.5,3);
\draw[blue,very thick] (-3,-1.5)--(3,1.5);
\draw[teal,densely dotted,thick] (0,0)--(1,0)--(1,2);
\draw[blue,densely dotted,thick] (0,0)--(0,1)--(2,1);
\node[teal,below=5pt] at (.5,0) {$1$};
\node[teal] at (1.4,1.9) {$2$};
\node[blue,left=5pt] at (-.2,.8) {$1$};
\node[blue] at (1.75,1.35) {$2$};
\fill[teal] (1,2) circle(2pt);
\fill[blue] (2,1) circle(2pt);

\end{tikzpicture}\par\smallskip
{\small \textcolor{teal}{\rule[.5ex]{1em}{.8pt}\ $y=2x,\ k=2/1=2$}\qquad\textcolor{blue}{\rule[.5ex]{1em}{.8pt}\ $y=x/2,\ k^\prime=1/2$}}\par
{\small 横纵增量互换，斜率互为倒数，符号保持不变。}\par\smallskip
\end{minipage}
\end{center}


取 $b=f(a)$，另一点 $(a+h,f(a+h))$，设 $k=f(a+h)-f(a)$（这里 $k$ 表示增量，不是斜率）。因 $f$ 一一对应，$h\ne0$ 时 $k\ne0$。
\[
\frac{f^{-1}(b+k)-f^{-1}(b)}{k}
=\frac{h}{k}=\frac{1}{[f(a+h)-f(a)]/h}.
\]
当 $k\to0$，由反函数在 $b$ 连续，$h\to0$。若 $f'(a)\ne0$，取极限得到
\[
\boxed{(f^{-1})'(b)=\frac1{f'(a)}}.
\]
适用条件：在 $a$ 附近的区间上可逆、原函数在 $a$ 可导且导数非零、反函数在 $b$ 连续。例如连续严格单调的函数满足所需的反函数连续性。若原导数为零，不能直接代公式。
\end{proofbox}
